La corda d'un violí fa 32~cm de llargària i vibra amb una freqüència fonamental de 196~Hz.

\figura[0.4]{fig1.png}

\begin{apartat}{1}
Expliqueu raonadament quina és la longitud d'ona del mode fonamental i digueu en quins punts de la corda hi ha els nodes i els ventres. Calculeu la velocitat de propagació de les ones que, per superposició, han generat l'ona estacionària de la corda.
\end{apartat}

\begin{apartat}{1}
Dibuixeu, de manera esquemàtica, el perfil de l'ona estacionària del tercer i del cinquè modes de vibració i calculeu-ne les freqüències.
\end{apartat}
